\documentclass[10pt,a4paper]{article}

\usepackage[utf8]{inputenc}
\usepackage[T1]{fontenc}
\usepackage[french]{babel}
\usepackage[margin=1.6cm]{geometry}
\usepackage{amsmath}
\usepackage{xcolor}
\usepackage{titlesec}
\usepackage{parskip}
\usepackage{array}

% --- Palette : une couleur par feature ---
\definecolor{c1}{RGB}{0,120,110}    % UA        - sarcelle
\definecolor{c2}{RGB}{30,110,170}   % H_ve      - bleu
\definecolor{c3}{RGB}{198,93,54}    % C         - corail
\definecolor{c4}{RGB}{176,137,20}   % A_solaire - ambre
\definecolor{c5}{RGB}{120,72,150}   % compacite - violet
\definecolor{anthracite}{RGB}{70,70,70}

\titlespacing{\section}{0pt}{5pt}{2pt}
\titleformat{\section}{\normalfont\large\bfseries}{}{0pt}{}

\newcommand{\feat}[2]{\section{\textcolor{#1}{#2}}}
\newcommand{\ou}[1]{{\small\textcolor{anthracite}{\textit{où :} #1}}\par}
\newcommand{\src}[1]{{\footnotesize\textcolor{anthracite}{\textbf{Sources :} \texttt{#1}}}\par}
\newcommand{\attr}[1]{\texttt{#1}}

\setlength{\parskip}{2pt}

\begin{document}

\begin{center}
{\LARGE\bfseries\textcolor{c1}{Modèle réduit d'enveloppe thermique} — \textcolor{c3}{5 features physiques}}
\end{center}

\vspace{-2pt}
{\small Plutôt qu'une quarantaine d'attributs bruts ResStock (catégories américaines : «~Wood Stud R-15~»,
zones ASHRAE\ldots), on agrège \textbf{cinq grandeurs physiques universelles} — transférables,
interprétables et directement exploitables par un modèle RC. Elles décrivent les déperditions, les
apports solaires, l'inertie et la forme du bâtiment.}

\vspace{2pt}
\noindent\colorbox{black!5}{\begin{minipage}{0.985\textwidth}
{\small\textbf{Notations générales :}
$A$ = surface d'une paroi (m\textsuperscript{2}) \;\textbullet\;
$R$ = résistance thermique (m\textsuperscript{2}K/W) \;\textbullet\;
$U = 1/R$ = conductance (W/m\textsuperscript{2}K) \;\textbullet\;
$V$ = volume chauffé (m\textsuperscript{3}) \;\textbullet\;
$A_{\text{pl}}$ = surface de plancher (m\textsuperscript{2}).}
\end{minipage}}

\feat{c1}{1.\quad UA — Déperdition par transmission (W/K)}
Chaleur traversant l'enveloppe (murs, toit, sol, fenêtres) par degré d'écart intérieur/extérieur.
Chaque paroi est un chemin de fuite ; on somme les conductances $U\!\cdot\!A$ (parois en parallèle).
Plus $UA$ est grand, plus le bâtiment fuit.
\[
UA = \tfrac{A_{\text{mur}}}{R_{\text{mur}}} + \tfrac{A_{\text{toit}}}{R_{\text{toit}}}
   + \tfrac{A_{\text{sol}}}{R_{\text{sol}}} + U_{\text{fen}}\,A_{\text{fen}}
\]
\ou{$R_{\text{mur/toit/sol}}$ et $U_{\text{fen}}$ = isolation de chaque paroi ; toiture routée selon la présence de combles.}
\src{in.insulation\_wall/roof/floor, in.window\_ufactor, out.params.*\_area, in.geometry\_attic\_type}

\feat{c2}{2.\quad H\textsubscript{ve} — Déperdition par renouvellement d'air (W/K)}
L'air neuf entrant par les fuites doit être réchauffé en continu : perte proportionnelle au débit
d'air et au volume.
\[
H_{ve} = 0{,}34 \; n \; V, \qquad n = \tfrac{\text{ACH50}}{20}, \qquad V = A_{\text{pl}} \times 2{,}5
\]
\ou{$0{,}34$ = chaleur d'un m\textsuperscript{3} d'air par °C ; $n$ = renouvellement réel (/h) déduit de la perméabilité mesurée ACH50 ; hauteur sous plafond 2,5~m.}
\src{in.air\_leakage\_to\_outside\_ach50, out.params.floor\_area\_lighting}

\feat{c3}{3.\quad C — Inertie thermique (kJ/K)}
Capacité à stocker la chaleur : un bâtiment lourd (béton, brique) amortit les variations, un léger
(bois) suit la météo. Faute de données de masse, estimée par classe ISO 13790.
\[
C = C_{\text{classe}} \times A_{\text{pl}}, \qquad
C_{\text{classe}} \in \{110,\,165,\,260,\,370\}\ \text{kJ/m}^2\text{K}
\]
\ou{$C_{\text{classe}}$ = capacité par m\textsuperscript{2} selon le mur : bois\,$\to$\,léger (110), acier\,$\to$\,moyen (165), brique\,$\to$\,lourd (260), béton\,$\to$\,très lourd (370).}
\src{in.geometry\_wall\_type, in.wall\_finish\_brick, out.params.floor\_area\_lighting}

\feat{c4}{4.\quad A\textsubscript{solaire} — Surface solaire équivalente (m\textsuperscript{2})}
Apports solaires gratuits par les vitrages. Équivalent d'un capteur parfait de même énergie, pondéré
par l'orientation (le sud capte le plus).
\[
A_{\text{solaire}} = \text{SHGC}\times 0{,}9 \times \sum_{\text{façades}} A_{\text{façade}}\; w
\]
\ou{SHGC = fraction de soleil transmise par le vitrage ; $0{,}9$ = facteur de cadre ; $w$ = poids d'orientation (Sud 1{,}0 ; E/O 0{,}6 ; Nord 0{,}3).}
\src{in.window\_shgc, in.orientation\_sin/cos, in.window\_areas, out.params.window\_area}

\feat{c5}{5.\quad Compacité — Coefficient de forme (m\textsuperscript{-1})}
À volume égal, une forme ramassée expose moins de surface qu'une forme étalée : elle perd moins.
Faible = compact (immeuble) ; élevée = dispersé (petite maison).
\[
\text{compacité} = \frac{\text{surface déperditive}}{V}
\]
\ou{surface déperditive = somme des parois donnant sur l'extérieur (murs, toit, sol, fenêtres, portes).}
\src{out.params.*\_area, volume estimé $A_{\text{pl}}\times 2{,}5$}

\vspace{3pt}
\noindent\colorbox{c1!8}{\begin{minipage}{0.985\textwidth}
\textbf{\textcolor{c1}{Exemple global}} — un logement du jeu de données (murs peu isolés, simple vitrage) :

\smallskip
{\small
\begin{tabular}{@{}>{\bfseries}l l | >{\bfseries}l l@{}}
\textcolor{c1}{UA} & $390$ W/K & \textcolor{c3}{C} & $11\,600$ kJ/K \\
\textcolor{c2}{H\textsubscript{ve}} & $75$ W/K & \textcolor{c4}{A\textsubscript{solaire}} & $4{,}4$ m\textsuperscript{2} \\
\multicolumn{2}{l|}{\textbf{Déperdition totale} $H = UA + H_{ve} = 465$ W/K} & \textcolor{c5}{compacité} & $0{,}85$ m\textsuperscript{-1} \\
\end{tabular}}

\smallskip
{\small\textit{Lecture :} par 20~°C dedans et 0~°C dehors ($\Delta T=20$), les pertes valent
$H\times\Delta T = 465\times20 \approx 9{,}3$~kW ; l'inertie $C$ stocke $\approx 3{,}2$~kWh par °C,
ce qui lisse les variations ; l'\textcolor{c4}{$A_{\text{solaire}}$} et la \textcolor{c5}{compacité}
modulent respectivement les apports gratuits et l'exposition de la forme.}
\end{minipage}}

\end{document}
